\documentclass[10pt]{beamer}
\usetheme{Madrid}
\usecolortheme{default}

\title{GSM A5/1 Stream Cipher Implementation\\and Known-Plaintext Attack Simulator}
\author{Anis Boulfoul \and Houssem Beniakhlef \and Khaled Guessaoui}
\institute{Wireless and Mobile Network Security}
\date{\today}

\begin{document}

\begin{frame}
  \titlepage
\end{frame}

\begin{frame}{Project Goal}
\begin{itemize}
  \item Implement the GSM A5/1 stream cipher in Python.
  \item Simulate encryption/decryption of GSM-like frames.
  \item Demonstrate a known-plaintext attack in an educational reduced keyspace setup.
  \item Deliver a reliable live demo for presentation day.
\end{itemize}
\end{frame}

\begin{frame}{A5/1 Cipher Design}
\begin{itemize}
  \item Three LFSRs:
  \begin{itemize}
    \item R1: 19 bits
    \item R2: 22 bits
    \item R3: 23 bits
  \end{itemize}
  \item Majority clocking controls irregular register updates.
  \item Initialization inputs:
  \begin{itemize}
    \item 64-bit session key
    \item 22-bit frame number
  \end{itemize}
  \item Output keystream is XORed with plaintext/ciphertext.
\end{itemize}
\end{frame}

\begin{frame}{What We Implemented}
\begin{itemize}
  \item \texttt{src/a51.py}: cipher core (LFSRs, clocking, keystream, encrypt/decrypt).
  \item \texttt{src/attack.py}: known-plaintext attack simulator.
  \item \texttt{src/demo.py}: one-command guided end-to-end demo.
  \item \texttt{src/cli.py}: command-line interface:
  \begin{itemize}
    \item \texttt{keystream}, \texttt{encrypt}, \texttt{decrypt}
    \item \texttt{attack}, \texttt{attack-sweep}, \texttt{attack-rate}, \texttt{demo}
  \end{itemize}
\end{itemize}
\end{frame}

\begin{frame}{Known-Plaintext Attack (Educational)}
\begin{itemize}
  \item Attacker knows a plaintext segment and matching ciphertext.
  \item Keystream segment is derived using XOR.
  \item Simulator brute-forces only unknown lower key bits (reduced keyspace).
  \item Returns candidate keys and attack runtime.
  \item Important: this is a pedagogical simulator, not a full real-world A5/1 break.
\end{itemize}
\end{frame}

\begin{frame}{Reproducible Demo Scenario}
\begin{itemize}
  \item Real key: \texttt{0x123456789ABCDEF0}
  \item Frame number: \texttt{0x13456}
  \item Known plaintext: \texttt{GSM-A51-DEMO-BLOCK}
  \item Unknown bits for attack: \texttt{12} (4096 key candidates)
  \item Generated ciphertext:
  \begin{center}
    \texttt{6b82759166cc706f9c78326bc4fc140bbe18}
  \end{center}
  \item Recovered candidate key:
  \begin{center}
    \texttt{0x123456789ABCDEF0}
  \end{center}
\end{itemize}
\end{frame}

\begin{frame}{Demo Flow (Live)}
\begin{enumerate}
  \item Show initial R1, R2, R3 register states.
  \item Generate and show keystream.
  \item Encrypt plaintext and show ciphertext (hex).
  \item Run known-plaintext attack and recover key candidate.
  \item Decrypt ciphertext and recover original message.
\end{enumerate}
\vspace{0.5em}
Main live command:
\begin{center}
  \texttt{python3 -m src.cli demo}
\end{center}
\end{frame}

\begin{frame}{Evaluation Results}
\begin{itemize}
  \item Unit/integration tests passing: \textbf{7/7}.
  \item Attack sweep shows fewer candidates with longer known plaintext.
  \item In our fixed scenario:
  \begin{itemize}
    \item 1 byte known plaintext: multiple candidates
    \item 2+ bytes: unique key recovery in tested runs
  \end{itemize}
  \item Attack-rate mode provides reproducible percentages over multiple trials.
\end{itemize}
\end{frame}

\begin{frame}{Team Contribution}
\begin{itemize}
  \item \textbf{Anis Boulfoul}: A5/1 implementation and integration.
  \item \textbf{Houssem Beniakhlef}: attack simulation logic and evaluation.
  \item \textbf{Khaled Guessaoui}: demo workflow and presentation preparation.
\end{itemize}
\end{frame}

\begin{frame}{Conclusion}
\begin{itemize}
  \item We built a complete, reproducible GSM A5/1 educational security demo.
  \item The project links cipher internals with an attack perspective.
  \item The live terminal demo is stable and presentation-ready.
\end{itemize}

\vspace{1em}
\centering
Thank you.
\end{frame}

\end{document}
